\notepage{15b}{Diagonalizing within each degenerate eigenspace}
The eigenvalue problem for $\Lambda$ can be solved independently
inside every eigenspace of $\Omega$. Since each restricted block
$\Lambda_a$ is Hermitian, choose an orthonormal eigenbasis for that
block. This replaces the old basis by linear combinations within
each block, with no mixing between distinct $\omega$ eigenspaces.

\smallheading{The block transformation}
Let $W_a$ diagonalize $\Lambda_a$, and combine the transformations as
$W=\diag(W_1,W_2,\ldots)$. Then
\[
 W^\dagger[\Lambda]W=
 \begin{bmatrix}
 D_1&0&\cdots\\0&D_2&\cdots\\\vdots&\vdots&\ddots
 \end{bmatrix},\qquad
 D_a=\diag(\lambda_{a1},\ldots,\lambda_{am_a}).
\]
The matrix of $\Omega$ is unchanged because every block is a multiple
of the identity:
\[
 W^\dagger[\Omega]W
 =\diag(\omega_1 I_{m_1},\omega_2 I_{m_2},\ldots)=[\Omega].
\]
\begin{center}
\begin{tikzpicture}[>=Stealth,node distance=1.1cm]
\node[draw=ink,rounded corners,align=center,text width=5.2cm,inner sep=8pt] (a)
 {Original $\Omega$ eigenbasis\\$\Omega$: diagonal\\$\Lambda$: block diagonal};
\node[draw=ink,rounded corners,align=center,text width=5.2cm,inner sep=8pt,right=1.5cm of a] (b)
 {Joint eigenbasis\\$\Omega$: same diagonal matrix\\$\Lambda$: diagonal};
\draw[->,thick,ink] (a)--node[above,font=\small]{$W$}(b);
\end{tikzpicture}
\end{center}
The diagram expresses the source's observation: $\Omega$ has the same
matrix in the two bases, while the remaining blocks of $\Lambda$
become diagonal.

\smallheading{Labeling simultaneous eigenvectors}
The new vectors are eigenvectors of both operators. A more informative
label therefore includes both eigenvalues:
\[
 \Omega\ket{\omega,\lambda}=\omega\ket{\omega,\lambda},
 \qquad
 \Lambda\ket{\omega,\lambda}=\lambda\ket{\omega,\lambda}.
\]
If these two eigenvalues still leave a degeneracy, an additional
degeneracy label can be retained. The construction proves the
simultaneous-diagonalization statement on note~15a.
